% commentaire.tex — f173v, vv. 920–958 (conversion initiale des travaux antérieurs)
\comm{922}{ᾐγυπτιάζετ᾽.}{Αἰγυπτιάζειν, faire l'Égyptien, comme Σιφνιάζειν, Χιάζειν, \italique{faire le Chiote}, Ἑλληνίζειν, \italique{faire l'Hellène}, Λεσβιάζειν, Βακίζειν, \italique{faire le Bakis} (\italique{Paix}, 1072). Les Égyptiens avaient la réputation d'être très rusés. On a conservé un vers d'Eschyle\nc{Tiré d'une scholie à Théocrite (15, 48b). C'est le fragment 373 de l'édition de Radt.} : \citComm{δεινοί πλέκειν τοι μηχανὰς Αἰγύπτιοι\\
Les Égyptiens habiles à ourdir des trames.}Cf. Properce (III, 9, 33) : \latin{Noxia Alexandria, dolis aptissima lellus.} (Note de Willems ; cf. aussi Taillardat §409.)}
\comm{928}{αὕτη μὲν ἡ μήρινθος οὐδὲν ἔσπασεν.}{Le tour est proverbial et l'expression nous est bien connue grâce à une scholie au vers 175 des \italique{Guêpes} où Aristophane joue avec l'halieutique proverbe :\citComm{παρὰ τὴν παροιμίαν « αὕτη μὲν ἡ μήρινθος οὐδὲν ἔσπασεν ». ἐλέγετο δὲ ἐπὶ τῶν πείρᾳ καθιέντων ἄγκιστρον μετὰ σχοίνου, καὶ μὴ ἐπιτυγχανόντων ἰχθύος τινός.\\
\francais{En regard du proverbe : « voici une ligne qui n'a rien tiré hors de l'eau ». On le dit à propos de ceux qui jettent un appât à la pointe acérée (un hameçon avec un roseau), mais qui ne rencontrent pas un poisson.}}Ce proverbe est dit d'une ruse qui échoue\nc{Pour un exposé plus développé sur ces deux passages avec à la clef l'explication du jargon halieutique, le lecteur pourra aller consulter le §395 des \italique{Images d'Aristophane} de J. Taillardat. Érasme a recensé l'adage (n° 845).}. En français : un échec sur toute la ligne.}
\comm{935}{ἱστιορράφος.}{Allusion aux hardes dont était couvert le prétendu Ménélas. (Note de Willems ; ἱστιορράφος est mis παρὰ προσδοκίαν pour μηχανορράφος\nc{Cf. Euripide, \italique{Andromaque}, vv. 447, 1116 ; Sophocle, \italique{Œdipe Roi}, v. 387.}.)}
\comm{936–7}{ὦ πρύτανι πρὸς τῆς δεξιᾶς, ἥνπερ φιλεῖς / κοίλην προτείνειν ἀργύριον ἤν τις διδῷ.}{Précis à l'usage du justiciable confronté à un prytane.
Les deux répliques que Mnésiloque adresse au prytane nous donnent un bon aperçu de l'estime en laquelle Aristophane tenait ces magistrats -- car nous allons voir que Mnésiloque tente de le soudoyer, comme cela semble l'usage à Athènes.
On trouve le même reproche de concussion dans la \italique{Paix} (vers 905-908) :\citComm{Ἀλλ᾽, ὦ πρυτάνεις, δέχεσθε τὴν Θεωρίαν.\\
Θέασ᾽ ὡς προθύμως ὁ πρύτανις παρεδέξατο.\\
Ἀλλ᾽ οὐκ ἄν, εἴ τι προῖκα προσαγαγεῖν σ᾽ ἔδει,\\
ἀλλ᾽ ηὗρον ἄν σ᾽ ὑπέχοντα τὴν ἐκεχειρίαν.\\
Allons, prytanes, recevez Théoria. Voyez l'accueil empressé que lui fait ce prytane. Quelle différence s'il t'avait fallu introduire une affaire gratis. Je t'aurais vu te retrancher sur les vacances.}(ὑπέχοντα τὴν ἐκεχειρίαν, «~alléguant la trêve~», est la version bégayée de ce qu'il fallait vraiment entendre : ὑπέχοντα τὴν χεῖρα, «~tendant la main~».)\\
En effet, pour soumettre une requête à la Boulé l'on devait préalablement obtenir l'aval des prytanes ; le cas échéant, les bouleutes discutaient et rédigeaient une proposition de loi (προβούλευμα) soumise ensuite à l'Ecclésia. L'auteur de la \italique{République des Athéniens} (III, 1) constate qu'il fallait parfois patienter jusqu'à une année avant que les prytanes ne donnent leur accord.
De même, pour intenter une action en justice, un citoyen devait régler des \italique{frais de justice}, τά πρυτανεῖα\nc{Voir à ce sujet le Daremberg \& Saglio à l'entrée \latin{Prytaneia}. Ces informations aident à interpréter la glose d'Hésychius (π 1511) : \citComm{πεπρυτανευκώς· ἐπιμεμηνιευκώς. καὶ προανηλωκώς,}rangée parmi les fragments comiques adespotes (408) dans l'édition Kassel-Austin.}, dont le nom indique clairement à qui cette somme devait être remise.\\ 
Par ses fonctions, le prytane \italique{s'engageait} à servir d'intermédiaire : ce genre de serment passait très certainement par une \italique{poignée de main}\nc{On pourra lire l'article de Francesco Mari : «~Les sens de la poignée de main en Grèce ancienne du VIIIe au Ve siècle avant J.-C.~», in \italique{Ktèma}, 43 (2018), pp. 105-131.}.
C'est le sens de «~ἥνπερ φιλεῖς / κοίλην προτείνειν~», «~[cette main] que tu as l'habitude de tendre ouverte~».\\
Mais toutes ces transactions étaient autant d'occasions pour les prytanes de détourner à leur profit une partie du bien public.
C'est entendu : le prytane est concussionnaire mais il serait suicidaire pour Mnésiloque de le lui dire !\citComm{Flattez-les, payez-les d'agréables mensonges,\\
Quelque indignation dont leur cœur soit rempli,\\
Ils goberont l'appât, vous serez leur ami.}Au contraire : Mnésiloque tente de mettre le prytane dans sa poche -- il y parvient d'ailleurs au vu de la bienveillante réponse qu'il reçoit : τί σοι χαρίσωμαι;\\
Pour cela, jeu de scène\nc{Il faut rendre aux jeux de scènes toute l'importance qu'ils avaient à l'époque ; empêtrés dans leurs costumes et d'imposants masques de théâtre, il est à parier que les acteurs compensaient cette gêne par des jeux de mains des plus expressifs.} : tout en déclamant «~πρὸς τῆς δεξιᾶς κτλ.~», Mnésiloque, qui n'est pas encore lié, tend sa main vers le prytane et s'offre par ce geste en suppliant\nc{(La \italique{droite} en question est bien celle de Mnésiloque et non du prytane.) Pour cette signification de la main tendue, cf. \LSJ \ à l'entrée προτείνω II.}. 
(Et d'ajouter -- en aparté, évidemment ! mais ce que le prytane a bien compris et qui explique ses bonnes dispositions à l'égard de Mnésiloque -- : «~ἀργύριον ἤν τις διδῷ~», «~[cette main que tu aimes à tendre] lorsque l'on t'offre de l'argent~».)\\
Le prytane accordera bien cette petite faveur à Mnésiloque (χάρισαι βραχύ τί μοι), à lui qui est sur le point de mourir ; et Mnésiloque d'insister sur ce dernier point (καίπερ ἀποθανουμένῳ).
Aristophane joue avec une locution proverbiale, selon sa pente maintenant bien établie par notre commentaire : les Athéniens, même à l'article de la mort, étaient réputés continuer de quémander la main tendue\nc{Recensé par Diogénianus (III, 12), \italique{Corpus paroemiographorum Graecorum}, I, p. 216. C'est l'adage n° 1646 d'Érasme. À noter qu'Aristophane y fait aussi une possible allusion aux vers 170 et suivants des \italique{Grenouilles}, lorsque Dionysos hèle un mort pour qu'il porte ses bagages, ce dernier lui répond : Δύο δραχμὰς μισθὸν τελεῖς; Voir plus haut le commentaire du vers 738.} :\citComm{Ἀττικὸς ὑπέχει τὴν χεῖρα ἀποθνήσκων· ἐπὶ τῶν φιλαργύρων. φιλοκερδεῖς γὰρ οἱ Ἀθηναῖοι.}Le \italique{verre du condamné} que demande Mnésiloque (καίπερ ἀποθανουμένῳ) n'est autre que cette faveur mendiée la main tendue vers le prytane.}
\comm{939–42}{γυμνὸν ἀποδύσαντά με / κέλευε πρὸς τῇ σανίδι δεῖν τὸν τοξότην, / ἵνα μὴ 'ν κροκωτοῖς καὶ μίτραις γέρων ἀνὴρ / γέλωτα παρέχω τοῖς κόραξιν ἑστιῶν.}{Ces quatre vers sont peut-être parmi les plus mécompris de toute la comédie.
L'erreur vient de ce que `πρός τῇ σανίδι δεῖν' signifie littéralement « lier \italique{contre} la planche » et non pas « lier à la planche » ; cette dernière expression se construisant en grec avec la préposition ἐν et non πρός\nc{Cf. les vers 934 (δῆσον αὐτὸν [...] ἐν τῇ σανίδι) et 1165 (οὑν [$=$ ὁ ἐν] τῇ σανίδι).}.
Que dit en apparence Mnésiloque ? il fait profession honorable de pudeur ; ce pudibond de la onzième heure tient à cacher ce qui doit être caché. 
Voilà pourquoi il demande à être retourné \italique{contre} (πρός) \italique{la planche} du supplice (τῇ σανίδι)\nc{En évitant par la même occasion que les corbeaux ne l'énucléent et ne se repaissent de ses yeux... spectacle qu'on image sans peine peu ragoûtant et fort désagréable pour le supplicié.}.
À qui tient-il à cacher son anatomie ? aux femmes évidemment\nc{Et non aux corbeaux, comme on peut le lire dans les éditions les plus récentes des \italique{Thesmophories}...} !.. Au vers 942, les `κροκωτοῖς καὶ μίτραις' sont les \italique{crocotes} et \italique{mitres} des Athéniennes réunies, rappelons-le, en assemblée dans le but de perdre Euripide. Dans ce contexte, ἐν + datif a le sens qu'il a toujours\nc{\italique{Confer} les vers 90-91.} : \italique{au milieu de}, \italique{parmi}, \italique{devant} : «~ἐν κροκωτοῖς καὶ μίτραις~» $=$ «~au milieu des (en)crocotées et des (en)mitrées~».
Voici pour le sens clair.

Mais n'oublions pas que Mnésiloque pratique avec un art consommé le double langage : toute parole qu'il adresse au prytane est à double sens : un clair, où il est dit que le sens véritable est caché\nc{Quelques vers auparavant, il avait demandé une faveur au prytane, tout en lui faisant bien comprendre que cette faveur mériterait salaire. Sur la concussion des prytanes, voir aussi la \italique{Paix} (vers 905-908).}. 
Que demande réellement Mnésiloque au prytane ?
L'expression `δῆσαι πρὸς σανίδι' a heureusement été relevée par Hésychius (δ 926) -- et c'est le seul des nombreux lexique consultés à avoir relevé ce tour :\citComm{δῆσαι πρὸς σανίδι· Ἀττικῶς· βασανίσαι ἐν τῷ δεσμωτηρίῳ.\\
Lier face à la planche : Atticisme. Interroger dans la prison.}Mnésiloque veut un procès équitable (il sent que la justice d'exception risque de lui être par trop expéditive) : il est préférable pour lui de s'en remettre aux masculines institutions. Il demande donc à être conduit en prison pour y être interrogé.

Réponse négative du prytane : c'est sur place, \italique{devant les personnes présentes} et non pas au δεσμωτήριον, qu'il devra expier son forfait :\citComm{Πρύτανις --- ἔχοντα ταῦτ᾽ ἔδοξε τῇ βουλῇ σε δεῖν,\\
ἵνα τοῖς παροῦσι δῆλος ᾖς πανοῦργος ὤν.\\}La σανίς de l'expression rapportée par Hésychius désigne sans doute la \italique{planche} sur laquelle était inscrite l'ordre du jour des tribunaux.
Cf. en \italique{Guêpes} (v. 848) et Photius (σ 68) : \citComm{σανίς· καὶ ἐν ᾗ τοὺς κακούργους ἔδουν καὶ ἐν ᾗ τὰς δίκας ἐνεχάραττον· τὸ µὲν ἐν Θεσµοφοριαζούσαις, τὸ δὲ ἐν Σφηξὶ (848) δεδήλωται.}Aristophane, ainsi qu'il aime à le faire assez souvent, détourne un mot ou une expression imagée qu'il interprète littéralement.

Devant le refus du prytane, Mnésiloque se retourne contre sa crocote -- au cœur des péripéties de la pièce -- qu'il accuse d'être à l'origine de la situation : il la gronde comme on gronderait un enfant\nc{Voir par exemple, \italique{Philoctète}, v. 940 : οἷ᾽ ἔργ᾽ ὁ παῖς μ᾽ ἔδρασεν οὑξ Ἀχιλλέως.} :\citComm{Μνησίλοχος --- ἰαππαπαιάξ· ὦ κροκώθ᾽ οἷ᾽ εἴργασαι.}
C'est encore Hésychius (κ 4170) qui nous permet de confirmer qu'il y a probablement là un jeu de mots κρόκαλ' $\sim$ κροκώθ' :\citComm{κρόκαλον· τὸ πανοῦργον παιδίον.}}
\comm{943–56.}{Le papyrus d'Oxyrhynque n° 3839\marqueur{1},}{dont il a été déjà question aux vers 742-766 contient contient également un fragment de ce passage. Il est si gâté qu'il ne nous apporte aucune information (si ce n'est une colométrie que Lætitia suivit \anglais{without strong conviction}\nc{L. P. E. Parker, \italique{The Songs of Aristophanes} (1997), p. 430.}). Le lecteur le trouvera néanmoins reproduit en annexe ; gage à lui d'en tirer quelque chose. \nc{\italique{The Oxyrhynchus Papyri}, vol. LVI (1989), n°3839, fr. 2. Analysé par W. E. H. Coockle avec le concours de Colin Austin.}}
\comm{945}{ὦ κροκώθ᾽ οἷ᾽ εἴργασαι·}{Y a-t-il un jeu de mots entre κρόκαλον et κροκωτός ? 
Cf. Hésychius (κ 4170 et 4270) :\citComm{κρόκαλον· τὸ πανοῦργον παιδίον.}}
\comm{947–1000}{}{Pour une analyse plus fine du chant mêlé de danses de tout ce passage, je renvoie à l'article d'Ester Cerbo\nc{«~Uno spettacolo di danza alla festa delle donne: a proposito di Ar. \italique{Thesm.} 947-1000~» (2021).}.
Quant à sa structure, celle d'Enger, assez communément admise -- et notamment par Lætitia\nc{Dans son indispensable \italique{The Songs of Aristophanes}, pp. 428-436.} --, me semble juste  :\citComm{--- anapestes (v. 947-952) ;\\
--- proode (v. 953-958) ;\\
--- première paire strophique (v. 959-962 $=$ v. 963-965) ;\\
--- première mésode (v. 966-968) ;\\
--- seconde paire strophique (v. 969-976 $=$ v. 977-984) ;\\
--- seconde mésode (v. 985-989) ;\\
--- tierce paire strophique (v. 990-994 $=$ v. 995-1000).}}
\comm{949–52}{Παύσων σέβεται καὶ νηστεύει κτλ.}{Pauson était connu pour son indigence. Les parémiographes nous ont rapporté à son sujet deux proverbes\nc{Respectivement Michel Apostolius (XIII, 2) et Suidas (π 823).} :\citComm{Παύσωνος πτωχότερος· οὗτος ζωγράφος ἦν καὶ ἐπὶ πενίᾳ διατεθρύλλητο.\\
Παύσων καὶ Ἶρος: ὀνόματα κύρια πενήτων. ὁ Ἀσκληπιὸς Παύσωνα καὶ Ἶρον κἂν ἄλλον τινὰ τῶν ἀπόρων ἰάσαιτο. κτλ.\\
(Iros, le mendiant odieux d'Homère, \italique{Odyssée} XVIII.)}L'on trouvera tous les témoignages sur Pauson d'Athènes dans le \italique{Recueil Milliet} d'Adolphe Reinach\nc{\italique{Recueil Milliet : Textes grecs et latins relatifs à l'histoire de la peinture ancienne}, tome I, publiés, traduits et commentés par Adolphe Reinach (1921) ; §§ 131-132 ; 173-179.}. Pour ce qui est d'Aristophane, la scholie à notre passage indique :\citComm{Παύσων ὁ ζωγράφος εἰς πενίαν διαβάλλεται, καὶ διὰ τοῦτο εἶπε `νηστεύει'.}Pauson l'importun est aussi mentionné au vers 853 des \italique{Acharniens} :\citComm{οὐδ᾽ αὖθις αὖ σε σκώψεται\\
Παύσων ὁ παμπόνηρος\\
Ni les brocards de Pauson, le gueux achevé, tu ne les essuieras [sur le marché].\\
Σ : οὗτος ὁ Παύσων ζωγράφος πένης σχωπτολόγος.}Enfin au vers 602 du \italique{Ploutos} :\citComm{Χρέμυλος --- Παρὰ τῆς Ἑκάτης ἔξεστιν τοῦτο πυθέσθαι,\\
εἴτε τὸ πλουτεῖν εἴτε τὸ πεινῆν βέλτιον. Φησὶ γὰρ αὕτη\\
τοὺς μὲν ἔχοντας καὶ πλουτοῦντας δεῖπνον κατὰ μῆν᾽ ἀποπέμπειν,\\
τοὺς δὲ πένητας τῶν ἀνθρώπων ἁρπάζειν πρὶν καταθεῖναι.\\
Ἀλλὰ φθείρου καὶ μὴ γρύξῃς\\
ἔτι μηδ᾽ ὁτιοῦν.\\
Οὐ γὰρ πείσεις, οὐδ᾽ ἢν πείσῃς.\\
Πενία --- Ὦ πόλις Ἄργους, κλύεθ᾽ οἷα λέγει.\\
Χρέμυλος --- \textbf{Παύσωνα} κάλει τὸν ξύσσιτον.\\
Chrémyle --- Il n'est que de s'informer chez Hécate, lequel vaut mieux d'être riche ou d'être pauvre. Car elle prétend que les riches et les opulents lui font hommage chaque mois d'un dîner, et que les pauvres s'en emparent avant qu'on l'ait mis là. Mais la peste soit de toi ; ne souffle plus le mot. Car tu ne me persuaderas pas, quand tu me persuaderais.\\
La Pauvreté --- « Cité des Argiens, écoutez ce qu'il dit ! »\\
Chrémyle --- Invoque \textbf{Pauson}, ton commensal.\\
Σ : ὁ Παύσων δὲ ἐπὶ πενίᾳ κωμῳδεῖται ζωγράφος ὦν. --- μετακαλοῦ σύντροφον τὸν Παύσωνα.}On constate que seuls les commentateurs indiquent que Pauson fût peintre (ζωγράφος) : il n'en n'est jamais question chez Aristophane.
Il paraît douteux que ce Pauson-là, le mendiant fâcheux à la dévotion fort intéressée\nc{Si sa dévotion prend les formes habituelles du souhait de renouvellement des prières (cf. le vers 287 ; Euripide, \italique{Électre}, vers 805), on se doute bien de ce qui attirait réellement Pauson-crève-la-faim. Voir ce que nous avons déjà dit aux vers 558-9 de la κρεανομία -- ou orgie institutionnalisée -- qui seule explique la réputation de piété des Athéniens.}, fût le peintre connu qu'évoque Aristote\nc{Aristote, \italique{Politique}, 1340a36 ; \italique{Poétique}, 1448a6.} ; la confusion venant peut-être de ce que la façon de peindre les figures du ζωγράφος ainsi que la seule anecdote nous ayant été rapportée à son sujet témoignent d'un caractère porté à la misanthropie\nc{Pauson le peintre peignait des figures pires que ce qu'elles étaient en réalité et Aristote le déconseille pour l'éducation des jeunes gens. On lui avait commandé le tableau d'un cheval se roulant dans la poussière ; l'artiste dessine un cheval courant et soulevant un nuage de poussière autour de lui. Il y travaillait quand survient l'amateur. Celui-ci récrimine : ce n'est pas le sujet commandé. Pauson fait retourner par son garçon le tableau sens dessus dessous ; on voit alors un cheval renversé, qui se roule dans la poussière.}.}
\comm{955}{χειρὶ σύναπτε χεῖρα.}{La danse que commande le chœur est probablement ce que Maurice Emmanuel\nc{\italique{La danse grecque antique} (1896), §319.} appelle une \italique{danse des mains jointes}. Il s'agit donc bien d'une danse individuelle et non en couple : chaque danseur joignant la main à son autre main.}
